% =========================================================================
%  Generador DDS con STM32H503CBT6 - Diseño y diagramas de flujo
%  Compilar:  pdflatex dds_diseno.tex   (dos veces, por las referencias)
%  Requiere doc/datos_escalera.dat (lo genera "make test").
% =========================================================================
\documentclass[11pt,a4paper]{article}
\usepackage[utf8]{inputenc}
\usepackage[T1]{fontenc}
\usepackage[spanish,es-noshorthands]{babel}
\usepackage{lmodern}
\usepackage[margin=2cm]{geometry}
\usepackage{amsmath,amssymb}
\usepackage{booktabs,array,tabularx}
\usepackage{xcolor}
\usepackage{tikz}
\usetikzlibrary{shapes.geometric,shapes.misc,arrows.meta,positioning,calc,fit,backgrounds}
\usepackage{pgfplots}
\pgfplotsset{compat=1.17}
\usepackage[hidelinks]{hyperref}
\usepackage{caption}

\definecolor{azul}{RGB}{31,78,140}
\definecolor{celeste}{RGB}{220,232,247}
\definecolor{verde}{RGB}{226,243,226}
\definecolor{amarillo}{RGB}{255,245,214}
\definecolor{rojo}{RGB}{250,224,224}
\definecolor{gris}{RGB}{238,238,238}

% ---------------- Estilos de diagrama de flujo ----------------
\tikzset{
  base/.style      = {draw, thick, align=center, font=\small, minimum height=8mm},
  inicio/.style    = {base, rounded rectangle, fill=verde, minimum width=26mm},
  proceso/.style   = {base, rectangle, fill=celeste, text width=58mm},
  procesoc/.style  = {base, rectangle, fill=celeste, text width=50mm},
  ancho/.style     = {base, rectangle, fill=celeste, text width=92mm},
  sub/.style       = {proceso, draw, double, double distance=1.2pt},
  decision/.style  = {base, diamond, aspect=2.4, fill=amarillo, inner sep=1pt, text width=30mm},
  io/.style        = {base, trapezium, trapezium left angle=72, trapezium right angle=108,
                      fill=rojo, text width=44mm, trapezium stretches body},
  flecha/.style    = {-{Stealth[length=2.6mm]}, thick},
  bloque/.style    = {draw, thick, rounded corners=2pt, align=center, font=\small,
                      fill=celeste, minimum height=10mm, inner sep=4pt},
  etq/.style       = {font=\footnotesize\itshape, inner sep=1pt},
}

\title{\textbf{Generador DDS y osciloscopio con STM32H503CBT6}\\[2mm]
\large Diseño, decisiones de ingeniería y diagramas de flujo}
\author{}
\date{}

\begin{document}
\maketitle
\vspace{-12mm}

\section{Resumen del diseño}

\begin{center}
\begin{tabular}{@{}ll@{}}
\toprule
\textbf{Parámetro} & \textbf{Valor} \\
\midrule
Microcontrolador & STM32H503CBT6 (Cortex-M33, 250\,MHz, 128\,KB flash, 32\,KB RAM) \\
Reloj de muestreo & TIM2 $\rightarrow$ TRGO, $F_s = 250\,\text{MHz}/250 = 1\,\text{MS/s}$ \\
Salidas & DAC1 CH1 (PA4) y CH2 (PA5), modo dual, 12 bits \\
Transferencia & GPDMA1 canal 0, circular por lista enlazada, buffer ping-pong $2\times256$ \\
Acumulador de fase & 32 bits $\Rightarrow$ resolución $F_s/2^{32} = 0{,}233$\,mHz \\
Tabla & 4096 muestras $\times$ 12 bits (en \texttt{uint16\_t}), doble buffer en RAM \\
Rango de salida & 0,1\,Hz a 400\,kHz ($0{,}4\,F_s$) \\
Interpolación & lineal, automática para $f < F_s/N = 244$\,Hz \\
Encoders & TIM1 (frecuencia) y TIM3 (forma de onda), modo encoder por hardware \\
Formas de onda & seno, triangular, sierra $\pm$, cuadrada, seno rectificado, arbitraria \\
\midrule
Osciloscopio & PA1 $\rightarrow$ ADC1 (único ADC del H503), 0 a 3,3\,V, disparado por TIM6 \\
Tiempo real & hasta 2\,MS/s a 12 bits; sobremuestreo por hardware hasta 16 bits \\
Tiempo equivalente & señales periódicas: 256 puntos/período, hasta $\approx 200$\,MS/s equivalentes \\
Salida de datos & USART2 a 2\,Mbaud por DMA $\rightarrow$ visor \texttt{tools/osciloscopio.py} \\
\bottomrule
\end{tabular}
\end{center}

% =========================================================================
\section{Arquitectura}

La idea central es que \textbf{el CPU nunca escribe en el DAC}. El timer marca
el ritmo con una precisión de cristal, el DAC convierte en cada flanco de TRGO
y el GPDMA le entrega la muestra siguiente. El CPU solo recalcula la mitad del
buffer que el DMA ya terminó de leer. Así el jitter de muestreo es nulo, sin
importar las interrupciones que haya.

\begin{figure}[h]
\centering
\begin{tikzpicture}[bloque/.append style={text width=22mm}]
  % columnas: x = 0 / 4.3 / 8.6 / 12.9 cm
  \node[bloque, fill=gris, text width=25mm] (encf) at (0,0)    {Encoder\\FRECUENCIA};
  \node[bloque, fill=gris, text width=25mm] (encw) at (0,-1.9) {Encoder\\FORMA};
  \node[bloque] (tim1) at (4.3,0)    {TIM1\\encoder x4};
  \node[bloque] (tim3) at (4.3,-1.9) {TIM3\\encoder x4};
  \node[bloque, fill=verde] (ui)  at (8.6,-0.95) {Lazo principal\\\texttt{ui\_task()}};
  \node[bloque, fill=gris, dashed] (lcd) at (8.6,1.3) {Pantalla\\(futuro)};
  \node[bloque, fill=amarillo] (par) at (12.9,-0.95) {Parámetros DDS\\$M$, tabla activa,\\modo CH2};
  \node[bloque, fill=rojo] (isr) at (12.9,-3.4) {ISR GPDMA\\\texttt{dds\_fill()}};
  \node[bloque] (buf)  at (12.9,-5.8) {Buffer RAM\\$2\times256\times32$\,b};
  \node[bloque] (dma)  at (8.6,-5.8)  {GPDMA1\\canal 0};
  \node[bloque] (dac)  at (4.3,-5.8)  {DAC1\\CH1 + CH2\\(DHR12RD)};
  \node[bloque] (tim2) at (4.3,-3.4)  {TIM2\\$F_s$ = 1\,MS/s};
  \node[bloque, fill=gris] (filt) at (0,-5.8) {Filtro\\reconstructor};

  \draw[flecha] (encf) -- (tim1);
  \draw[flecha] (encw) -- (tim3);
  \draw[flecha] (tim1.east) -- ++(6mm,0) |- ([yshift=2mm]ui.west);
  \draw[flecha] (tim3.east) -- ++(6mm,0) |- ([yshift=-2mm]ui.west);
  \draw[flecha] (ui) -- (par);
  \draw[flecha, dashed] (ui) -- (lcd);
  \draw[flecha] (par) -- (isr);
  \draw[flecha] (isr) -- node[etq, right, align=left]{escribe la\\mitad libre} (buf);
  \draw[flecha] (buf) -- (dma);
  \draw[flecha] ([yshift=-2.5mm]dma.west) -- node[etq, below]{datos} ([yshift=-2.5mm]dac.east);
  \draw[flecha] ([yshift=2.5mm]dac.east) -- node[etq, above]{req.} ([yshift=2.5mm]dma.west);
  \draw[flecha] (tim2) -- node[etq, right]{TRGO} (dac);
  \draw[flecha] (dma.north) |- node[etq, pos=0.75, above]{HT / TC} (isr.west);
  \draw[flecha] (dac) -- node[etq, above]{PA4/5} (filt);
\end{tikzpicture}
\caption{Diagrama en bloques del generador.}
\end{figure}

% =========================================================================
\section{Decisiones de diseño}

\subsection{¿TIM1 o TIM2?}

En el STM32H503 el DAC puede dispararse desde TIM1, TIM2, TIM3, TIM6, TIM7 y
LPTIM1/2 (campo \texttt{TSELx} de \texttt{DAC\_CR}). Todos los timers corren a
250\,MHz con el prescaler de APB en 1, así que la resolución de $F_s$ es la
misma con cualquiera. La diferencia está en lo que se pierde al usarlo:

\begin{center}
\begin{tabularx}{\linewidth}{@{}l>{\raggedright\arraybackslash}X>{\raggedright\arraybackslash}X@{}}
\toprule
\textbf{Timer} & \textbf{Características} & \textbf{Uso en este diseño} \\
\midrule
TIM1 & Avanzado, 16 bits, PWM complementario con tiempo muerto, el único del chip &
       \textbf{Encoder de frecuencia} (PA8/PA9). Usarlo solo como reloj desperdiciaría sus funciones. \\
TIM2 & 32 bits, 4 canales, dispara al DAC (\texttt{TSEL}=2) &
       \textbf{Reloj de muestreo del DDS} (elegido) \\
TIM3 & 16 bits, 4 canales, modo encoder & \textbf{Encoder de forma de onda} (PA6/PA7) \\
TIM6 & Básico, sin pines, pensado para el DAC (\texttt{TSEL}=5) &
       Alternativa equivalente: \texttt{DDS\_TIMER\_TIM6} en \texttt{dds\_config.h} libera TIM2 \\
\bottomrule
\end{tabularx}
\end{center}

Para generar el TRGO, TIM1 y TIM2 dan el mismo resultado. Elegí \textbf{TIM2}
porque TIM1 es el único timer avanzado del chip y su par de pines PA8/PA9 sirve
muy bien para el encoder. Hacen falta dos timers con modo encoder (TIM1 y TIM3),
y con TIM2 en el DDS no se desperdicia ninguna función especial.

\subsection{Frecuencia de muestreo y frecuencia máxima}

\begin{equation}
F_s = \frac{f_{TIM}}{ARR+1} = \frac{250\,\text{MHz}}{250} = 1\,\text{MS/s}
\qquad
M = \operatorname{round}\!\left(\frac{f_{out}\cdot 2^{32}}{F_s}\right)
\qquad
\Delta f = \frac{F_s}{2^{32}} = 0{,}233\,\text{mHz}
\end{equation}

Elegí 1\,MS/s porque es la velocidad de actualización típica que especifica ST
para sus DAC con el buffer de salida activado (confirmalo en la tabla
\emph{DAC characteristics} del DS14053). El límite teórico es Nyquist,
$F_s/2 = 500$\,kHz. En la práctica conviene $f_{max} = 0{,}4\,F_s = 400$\,kHz,
porque la primera imagen cae en $F_s - f$ y el filtro reconstructor necesita
banda de transición. La retención de orden cero atenúa según
$\operatorname{sinc}(f/F_s)$: a 400\,kHz se pierden 2,4\,dB. Para un seno limpio
sin un filtro muy abrupto, conviene quedarse por debajo de $F_s/4 = 250$\,kHz.
Las formas con armónicos (cuadrada, sierra) se ven bien hasta $\approx F_s/20$.

Si el DAC de tu placa lo soporta, \texttt{DDS\_FS\_HZ} acepta 2; 2,5 y 5\,MS/s,
que son divisores exactos de 250\,MHz. Por encima de 1\,MS/s usá
\texttt{MODE = 010} (buffer interno OFF) y un amplificador operacional rápido
como buffer externo. El CPU acompaña: a 5\,MS/s el lazo directo usa
$\approx 25$\,\% del tiempo.

\subsection{Tamaño de tabla recomendado: \texorpdfstring{$4096\times12$}{4096x12} bits}

En un DDS, $N$ \textbf{no} es la cantidad de muestras por período que salen. Esas
son siempre $F_s/f_{out}$. $N$ fija la \emph{precisión de fase}. Al truncar la
fase a $\log_2 N$ bits aparecen espurias de aproximadamente
$\text{SFDR}\approx 6{,}02\cdot\log_2 N$ dBc.

\begin{center}
\begin{tabular}{@{}rccccl@{}}
\toprule
$N$ & RAM (2 buffers) & SFDR trunc. & Error seno interp. & Sierra & Comentario \\
\midrule
256   & 1\,KB  & $\approx 48$ dBc & 0,15 LSB    & 16 LSB/paso & solo con interpolación \\
1024  & 4\,KB  & $\approx 60$ dBc & 0,01 LSB    & 4 LSB/paso  & aceptable \\
\textbf{4096} & \textbf{16\,KB} & $\boldsymbol{\approx 72}$ \textbf{dBc} & \textbf{0,0006 LSB} & \textbf{1 LSB/paso} & \textbf{recomendado} \\
16384 & 64\,KB & $\approx 84$ dBc & --          & 1 LSB/paso  & no entra en 32\,KB \\
\bottomrule
\end{tabular}
\end{center}

\textbf{Recomendación: 4096 muestras de 12 bits guardadas en \texttt{uint16\_t}}
(8\,KB por tabla), por estas razones:
\begin{itemize}\itemsep0pt
  \item 72\,dBc de SFDR coinciden con el techo de un DAC ideal de 12 bits
        ($6{,}02\cdot12+1{,}76 = 74$\,dB). Una tabla más grande no mejora la
        salida porque manda el DAC.
  \item Una sierra de 4096 pasos recorre los 4096 códigos: 1 LSB por muestra, sin pérdida.
  \item El índice se obtiene con un desplazamiento (\texttt{fase >> 20}), sin módulo.
  \item 12 bits justificados a la derecha van directo al registro \texttt{DHR12RD}.
  \item Con doble buffer (16\,KB) el cambio de forma de onda no produce glitches,
        y todavía sobra la mitad de la RAM.
  \item Se agrega una muestra guarda ($T[4096]=T[0]$) para que la interpolación
        lea $T[i+1]$ sin máscara.
\end{itemize}
Las tablas se calculan en el arranque (seno, triangular, etc.). Las formas
arbitrarias se ``graban'' en flash con \texttt{tools/gen\_tabla.py}, a partir de
una fórmula o de un CSV medido, y se copian a RAM al seleccionarlas.

\subsection{Interpolación lineal en baja frecuencia}

Cuando $f_{out} < F_s/N$ (244\,Hz), el paso de fase es menor que una muestra de
tabla y cada valor se repite varias veces. El resultado es una escalera. En ese
caso se usan los 12 bits de fracción que siguen al índice:

\begin{equation}
\underbrace{\texttt{fase}[31{:}20]}_{i}\;\;
\underbrace{\texttt{fase}[19{:}8]}_{\phi}\qquad
y = T[i] + \left\lfloor\frac{(T[i+1]-T[i])\cdot\phi + 2^{11}}{2^{12}}\right\rfloor
\end{equation}

El error máximo de la interpolación lineal sobre un seno es
$A(2\pi/N)^2/8 = 0{,}0006$\,LSB, despreciable. Medido en el test con el código
real: \textbf{0,99\,LSB máximo a 10\,Hz} (es el redondeo de la tabla más el de
la salida), contra 2,96\,LSB del camino directo.

La cuadrada y las sierras \textbf{no} se interpolan. Sus flancos se convertirían
en rampas de $1/N$ del período (244\,$\mu$s a 1\,Hz). Por encima de 244\,Hz se
usa el camino directo, más rápido, porque el error de truncamiento ya queda al
nivel del ruido del DAC. El umbral se cambia con \texttt{DDS\_INTERP\_MAX\_STEP}.

\begin{figure}[h]
\centering
\begin{tikzpicture}
\begin{axis}[width=0.9\linewidth, height=6.2cm, xlabel={muestra $n$ ($\mu$s)},
             ylabel={código DAC}, grid=major,
             legend pos=south east,
             legend style={font=\small}, tick label style={font=\small},
             /pgf/number format/use comma, /pgf/number format/1000 sep={}]
  \addplot[very thick, black!45] table[x=n, y=ideal] {datos_escalera.dat};
  \addlegendentry{seno ideal}
  \addplot[thick, red, const plot] table[x=n, y=sin_interp] {datos_escalera.dat};
  \addlegendentry{sin interpolación (escalera)}
  \addplot[thick, azul, const plot] table[x=n, y=con_interp] {datos_escalera.dat};
  \addlegendentry{con interpolación (firmware)}
\end{axis}
\end{tikzpicture}
\caption{Seno de 10\,Hz cerca del cruce por cero, salida real de \texttt{dds\_fill()}.
Sin interpolar, cada muestra de tabla se repite 24 veces (saltos de $\approx 3$\,LSB).
Interpolando, la salida avanza de a 1\,LSB.}
\end{figure}

\subsection{Presupuesto de CPU y memoria}

\begin{center}
\begin{tabular}{@{}lcccc@{}}
\toprule
\textbf{Camino (CH2 en 90°)} & \textbf{Instr./muestra} & \textbf{Ciclos (est.)} & \textbf{CPU 1\,MS/s} & \textbf{CPU 5\,MS/s} \\
\midrule
Directo     & 10 & $\approx 12$ & $\approx 5$\,\%  & $\approx 25$\,\% \\
Interpolado & 22 & $\approx 26$ & $\approx 10$\,\% & $\approx 52$\,\% \\
\bottomrule
\end{tabular}
\end{center}
A 1\,MS/s hay 250 ciclos por muestra. El firmware mide el peor caso real con el
contador DWT y lo muestra por UART (\texttt{CPU x.x\%}). Con el osciloscopio
incluido la RAM queda así: 16\,KB de tablas DDS + 2\,KB de buffer del DAC +
4\,KB de captura + 2,5\,KB del ETS + pila $\approx$ 28,5\,KB de 32\,KB.
Flash: 28,4\,KB de 128\,KB.

% =========================================================================
\clearpage
\section{Diagramas de flujo}

\begin{figure}[h!]
\centering
\begin{tikzpicture}[node distance=6mm]
  \node[inicio] (a) {Inicio (\texttt{main})};
  \node[sub, below=of a] (b) {\texttt{board\_init()}\\VOS0, flash 5 WS, PLL1 = 250\,MHz,\\ICACHE, SysTick 1\,ms, USART2};
  \node[sub, below=of b] (c) {\texttt{dds\_core\_init()}\\genera la tabla del seno (4097 muestras)\\$M = \text{round}(1\,\text{kHz}\cdot2^{32}/F_s)$};
  \node[sub, below=of c] (d) {\texttt{encoder\_init()}\\TIM1 y TIM3 en modo encoder x4,\\filtro de entrada, pull-ups, pulsadores};
  \node[sub, below=of d] (e) {\texttt{ui\_init()}\\estado inicial, \texttt{display\_init()} (weak)};
  \node[sub, below=of e] (f) {\texttt{dds\_hw\_start()}\\precarga el buffer, TIM2 (TRGO),\\DAC dual (HFSEL), GPDMA circular, CEN};
  \node[sub, below=of f] (g) {\texttt{ui\_task()}\\(ver Fig.~\ref{fig:ui})};
  \node[proceso, below=of g] (h) {\texttt{\_\_WFI()}: dormir hasta la próxima IRQ\\(DMA cada 256\,$\mu$s o SysTick)};
  \draw[flecha] (a) -- (b);
  \draw[flecha] (b) -- (c);
  \draw[flecha] (c) -- (d);
  \draw[flecha] (d) -- (e);
  \draw[flecha] (e) -- (f);
  \draw[flecha] (f) -- (g);
  \draw[flecha] (g) -- (h);
  \draw[flecha] (h.south) -- ++(0,-5mm) -| node[etq, pos=0.75, left]{lazo infinito} ($(g.west)+(-8mm,0)$) -- (g.west);
\end{tikzpicture}
\caption{Programa principal.}
\label{fig:main}
\end{figure}

\begin{figure}[p]
\centering
\begin{tikzpicture}[node distance=6mm and 12mm]
  \node[inicio] (s) {\texttt{ui\_task()}};
  \node[io, below=of s] (r1) {Leer clicks del encoder F\\$\Delta = (\text{CNT}_{\text{TIM1}} - \text{ant})/4$};
  \node[decision, below=of r1] (d1) {$\Delta \neq 0$?};
  \node[procesoc, right=of d1] (p1) {$f \leftarrow f + \Delta\cdot\text{paso}$, limitado a $[0{,}1\,\text{Hz};\,0{,}4F_s]$\\$M \leftarrow \text{round}(f\,2^{32}/F_s)$};
  \node[decision, below=12mm of d1] (d2) {¿Pulsador F?};
  \node[procesoc, right=of d2] (p2) {paso $\leftarrow$ siguiente\\0,1 / 1 / 10 / \dots / 100k Hz};
  \node[io, below=12mm of d2] (r2) {Leer clicks del encoder W\\(TIM3)};
  \node[decision, below=of r2] (d3) {$\Delta \neq 0$?};
  \node[procesoc, right=of d3] (p3) {onda $\leftarrow$ (onda $+\Delta$) mod 7\\generar tabla en el buffer \emph{inactivo}\\e intercambiar puntero (atómico)};
  \node[decision, below=12mm of d3] (d4) {¿Pulsador W?};
  \node[procesoc, right=of d4] (p4) {modo CH2 $\leftarrow$ siguiente\\90° / invertida / sync / igual};
  \node[decision, below=12mm of d4] (d5) {¿Hubo cambios?};
  \node[procesoc, right=of d5] (p5) {\texttt{display\_update()}\\+ línea por UART};
  \node[inicio, below=12mm of d5] (e) {Fin};

  \foreach \d/\n [count=\k] in {d1/d2, d2/r2, d3/d4, d4/d5, d5/e}
    \coordinate (m\k) at ($(\d.south)!0.5!(\n.north)$);

  \draw[flecha] (s) -- (r1);
  \draw[flecha] (r1) -- (d1);
  \draw[flecha] (r2) -- (d3);
  \foreach \d/\p/\n [count=\k] in {d1/p1/d2, d2/p2/r2, d3/p3/d4, d4/p4/d5, d5/p5/e} {
    \draw[flecha] (\d) -- node[etq, above]{sí} (\p);
    \draw[flecha] (\d) -- node[etq, right, pos=0.3]{no} (\n);
    \draw[flecha] (\p.south) |- (m\k);
  }
\end{tikzpicture}
\caption{Lazo de interfaz: los encoders cuentan por hardware y \texttt{ui\_task()}
solo lee la diferencia. Los cambios de $M$ entran en el borde del bloque
siguiente, con fase continua.}
\label{fig:ui}
\end{figure}

\begin{figure}[p]
\centering
\begin{tikzpicture}[node distance=6mm and 7mm, procesoc/.append style={text width=37mm}]
  \node[inicio] (s) {IRQ GPDMA1 canal 0};
  \node[decision, below=of s] (d0) {¿HTF\\(media transf.)?};
  \node[procesoc, right=11mm of d0] (h) {limpiar HTF\\destino = 1.ª mitad};
  \node[procesoc, left=11mm of d0] (t) {limpiar TCF\\destino = 2.ª mitad};
  \node[ancho, below=14mm of d0] (p) {Leer parámetros \textbf{una sola vez}:\\tabla activa $T$, palabra $M$, modo CH2, fase};
  \node[decision, below=of p, text width=34mm] (d1) {¿forma interpolable\\y $M < 2^{20}$?\\($f < F_s/N$)};
  \node[procesoc, left=of d1, fill=verde] (mi) {modo \textbf{interpolado}};
  \node[procesoc, right=of d1, fill=verde] (md) {modo \textbf{directo}};
  \node[ancho, below=14mm of d1] (i) {$i \leftarrow \text{fase} \gg 20$ \quad (índice de 12 bits)};
  \node[decision, below=of i] (d2) {¿interpolado?};
  \node[procesoc, left=of d2, text width=50mm] (yi) {$\phi \leftarrow (\text{fase}\gg8)\,\&\,\text{0xFFF}$\\$y_1 \leftarrow T[i] + (\Delta T\,\phi + 2^{11})\gg 12$\\$\Delta T = T[i+1]-T[i]$};
  \node[procesoc, right=of d2] (yd) {$y_1 \leftarrow T[i]$};
  \node[ancho, below=16mm of d2] (c2) {CH2 según el modo:\\
     90°: $y_2 = T[(\text{fase}+2^{30})\gg20]$ \qquad invertida: $y_2 = 4095-y_1$\\
     sync: $y_2 = 4095$ si fase $< 2^{31}$, si no 0 \qquad igual: $y_2 = y_1$};
  \node[ancho, below=of c2] (w) {\texttt{buf[k]} $\leftarrow y_1 \,|\, (y_2 \ll 16)$ \quad (formato DHR12RD)\\
     fase $\leftarrow$ fase $+ M$ \quad (desborde = módulo $2^{32}$)};
  \node[decision, below=of w] (d3) {¿$k < 256$?};
  \node[ancho, below=8mm of d3] (g) {guardar la fase y medir ciclos (DWT)};
  \node[inicio, below=of g] (e) {Fin de la ISR};

  \coordinate (m0) at ($(d0.south)!0.5!(p.north)$);
  \coordinate (m1) at ($(d1.south)!0.5!(i.north)$);
  \coordinate (m2) at ($(d2.south)!0.5!(c2.north)$);

  \draw[flecha] (s) -- (d0);
  \draw[flecha] (d0) -- node[etq, above]{sí} (h);
  \draw[flecha] (d0) -- node[etq, above]{no (TCF)} (t);
  \draw[thick] (h.south) |- (m0) -| (t.south);
  \draw[flecha] (m0) -- (p);
  \draw[flecha] (p) -- (d1);
  \draw[flecha] (d1) -- node[etq, above]{sí} (mi);
  \draw[flecha] (d1) -- node[etq, above]{no} (md);
  \draw[thick] (mi.south) |- (m1) -| (md.south);
  \draw[flecha] (m1) -- (i);
  \draw[flecha] (i) -- (d2);
  \draw[flecha] (d2) -- node[etq, above]{sí} (yi);
  \draw[flecha] (d2) -- node[etq, above]{no} (yd);
  \draw[thick] (yi.south) |- (m2) -| (yd.south);
  \draw[flecha] (m2) -- (c2);
  \draw[flecha] (c2) -- (w);
  \draw[flecha] (w) -- (d3);
  \draw[flecha] (d3) -- node[etq, right]{no} (g);
  \draw[flecha] (d3.east) -- node[etq, above]{sí, $k$++} ($(d3.east -| md.east)+(4mm,0)$) |- (i.east);
  \draw[flecha] (g) -- (e);
\end{tikzpicture}
\caption{Motor DDS (ISR del GPDMA). En el código, la decisión ``¿interpolado?''
y el modo de CH2 se resuelven \emph{antes} del lazo: el compilador genera 8
lazos especializados sin ningún \texttt{if} adentro.}
\label{fig:isr}
\end{figure}

\begin{figure}[h!]
\centering
\begin{tikzpicture}[node distance=5mm and 10mm]
  \node[inicio] (s) {SysTick (1\,ms)};
  \node[proceso, below=of s] (a) {\texttt{ms++}; para cada pulsador:\\leer pin (\texttt{IDR})};
  \node[decision, below=of a] (d1) {¿lectura $\neq$\\estado estable?};
  \node[procesoc, right=of d1] (r) {contador $\leftarrow$ 0};
  \node[proceso, below=8mm of d1] (c) {contador++};
  \node[decision, below=of c] (d2) {¿contador\\$\geq$ 20?};
  \node[proceso, below=8mm of d2] (u) {estado $\leftarrow$ lectura; contador $\leftarrow 0$\\si quedó apretado: evento = 1};
  \node[inicio, below=of u] (e) {Fin};
  \draw[flecha] (s) -- (a);
  \draw[flecha] (a) -- (d1);
  \draw[flecha] (d1) -- node[etq, above]{no} (r);
  \draw[flecha] (d1) -- node[etq, right]{sí} (c);
  \draw[flecha] (c) -- (d2);
  \draw[flecha] (d2) -- node[etq, right]{sí} (u);
  \draw[flecha] (d2.east) -- node[etq, above]{no} ++(15mm,0) |- (e.east);
  \draw[flecha] (r.south) |- (e.east);
  \draw[flecha] (u) -- (e);
\end{tikzpicture}
\caption{Antirrebote de los pulsadores por integración (20\,ms estables).}
\label{fig:deb}
\end{figure}

% =========================================================================
\clearpage
\section{Osciloscopio de 0 a 3,3\,V}

\subsection{¿ADC en paralelo y desfasados? El H503 tiene uno solo}

La idea de poner varios ADC en paralelo sobre la misma entrada, cada uno
disparado con un desfase de $T_s/N$ (\emph{time-interleaving}), multiplica por
$N$ la frecuencia de muestreo: mejora la \textbf{resolución en el tiempo}. Pero
el STM32H503 tiene \textbf{un único ADC} (ADC1). Lo confirman el header CMSIS
oficial, donde solo existe \texttt{ADC1}, y el datasheet: un ADC de 12 bits a
2,5\,MS/s. El modo intercalado necesita, por ejemplo, el STM32H563, que tiene
dos ADC.

Con un solo ADC se aprovecha el mismo principio de dos maneras:

\begin{center}
\begin{tabularx}{\linewidth}{@{}>{\raggedright\arraybackslash}p{33mm}>{\raggedright\arraybackslash}X>{\raggedright\arraybackslash}X>{\raggedright\arraybackslash}p{22mm}@{}}
\toprule
\textbf{Técnica} & \textbf{Qué mejora} & \textbf{Requisito} & \textbf{En el H503} \\
\midrule
$N$ ADC intercalados (desfase $T_s/N$) & Resolución en el tiempo $\times N$, con cualquier señal (incluso un evento único) & $N$ ADC sincronizados & No (1 ADC) \\
\textbf{Tiempo equivalente} (ETS): un ADC, muestras en fases distintas de períodos sucesivos & Resolución en el tiempo: 256 puntos por período, hasta $\approx 200$\,MS/s equivalentes & Señal \textbf{periódica} y estable durante la captura (1\,ms) & Sí, modo \texttt{ETS} \\
\textbf{Sobremuestreo}: promediar $2^n$ conversiones por muestra & Resolución en \textbf{amplitud}: $+\frac{1}{2}$ bit por cada $\times 2$, hasta 16 bits & Señal lenta frente a la velocidad del ADC & Sí, automático en bases $\leq 500$\,kS/s \\
\bottomrule
\end{tabularx}
\end{center}

\subsection{Arquitectura}

\begin{figure}[h]
\centering
\begin{tikzpicture}[bloque/.append style={text width=23mm}]
  \node[bloque, fill=gris] (in)  at (0,0)    {Entrada\\PA1\\0 -- 3,3\,V};
  \node[bloque, fill=gris] (buf) at (3.4,0)  {Buffer\\op-amp\\(externo)};
  \node[bloque] (adc) at (6.8,0)  {ADC1\\12 bits\\sobremuestreo};
  \node[bloque] (tim) at (6.8,2.2) {TIM6\\$F_s$ (TRGO)};
  \node[bloque] (dma) at (10.2,0) {GPDMA1\\canal 1};
  \node[bloque] (cap) at (13.4,0) {Captura RAM\\$2048\times16$\,b};
  \node[bloque, fill=verde, text width=30mm] (cpu) at (13.4,-2.6)
       {\texttt{osc\_task()}\\trigger, mediciones,\\tiempo equivalente};
  \node[bloque] (udma) at (9.0,-2.6) {GPDMA1\\canal 2};
  \node[bloque, fill=gris, text width=30mm] (pc) at (4.6,-2.6) {USART2 2\,Mbaud\\$\rightarrow$ PC\\\texttt{osciloscopio.py}};
  \node[bloque, fill=gris, dashed, text width=30mm] (lcd) at (13.4,-4.7) {Pantalla\\(\texttt{display\_scope})};

  \draw[flecha] (in) -- (buf);
  \draw[flecha] (buf) -- (adc);
  \draw[flecha] (tim) -- node[etq, right]{dispara} (adc);
  \draw[flecha] (adc) -- node[etq, above]{DR} (dma);
  \draw[flecha] (dma) -- (cap);
  \draw[flecha] (cap) -- node[etq, right]{TC} (cpu);
  \draw[flecha] (cpu) -- node[etq, above]{trama} (udma);
  \draw[flecha] (udma) -- (pc);
  \draw[flecha, dashed] (cpu) -- (lcd);
\end{tikzpicture}
\caption{Cadena de adquisición. Corre en paralelo con el DDS, que sigue
funcionando solo con TIM2, el DAC y el canal 0 del GPDMA.}
\end{figure}

El reloj del ADC es $\text{HCLK}/8 = 31{,}25$\,MHz. Una conversión de 12 bits
dura $2{,}5 + 12{,}5 = 15$ ciclos (480\,ns), así que el máximo en tiempo real
es \textbf{2\,MS/s}. Para llegar a los 2,5\,MS/s del datasheet hacen falta
37,5\,MHz, que se pueden obtener con el PLL2. Igual que en el DDS, el timer fija
el instante de cada conversión sin jitter y el DMA guarda los datos: el CPU solo
procesa cuando la captura terminó.

\subsection{Base de tiempo y sobremuestreo automático}

Para cada $F_s$ el firmware calcula PSC/ARR del timer y elige, en este orden:
(1)~el mayor sobremuestreo $R = 2^n$ (hasta 256) que entra en el período de
muestreo y (2)~con el tiempo que sobra, el mayor tiempo de muestreo del ADC, que
exige menos a la impedancia de la fuente. Todas las muestras se normalizan a 16
bits ($0 \ldots 65535 = 0 \ldots 3{,}3$\,V), así el formato no cambia.

\begin{center}
\begin{tabular}{@{}rrrcrc@{}}
\toprule
$F_s$ [S/s] & PSC & ARR & Sobremuestreo & SMP [ciclos] & Resolución efectiva \\
\midrule
2\,000\,000 & 0 & 124   & $\times1$   & 2,5  & 12,0 bits \\
1\,000\,000 & 0 & 249   & $\times1$   & 12,5 & 12,0 bits \\
500\,000    & 0 & 499   & $\times2$   & 12,5 & 12,5 bits \\
200\,000    & 0 & 1249  & $\times8$   & 6,5  & 13,5 bits \\
100\,000    & 0 & 2499  & $\times16$  & 6,5  & 14,0 bits \\
50\,000     & 0 & 4999  & $\times32$  & 6,5  & 14,5 bits \\
20\,000     & 0 & 12499 & $\times64$  & 6,5  & 15,0 bits \\
10\,000     & 0 & 24999 & $\times128$ & 6,5  & 15,5 bits \\
5\,000      & 0 & 49999 & $\times256$ & 6,5  & 16,0 bits \\
2\,000      & 1 & 62499 & $\times256$ & 47,5 & 16,0 bits \\
1\,000      & 3 & 62499 & $\times256$ & 92,5 & 16,0 bits \\
\bottomrule
\end{tabular}
\end{center}
La resolución efectiva es la teórica, $12 + \frac{1}{2}\log_2 R$. Se alcanza
cuando la señal trae al menos $\approx 1$\,LSB de ruido, que actúa como
\emph{dither}.

\subsection{Tiempo equivalente con un solo ADC}

Si la señal es periódica, no hace falta que varios ADC muestreen el mismo
período en instantes desfasados. Alcanza con que \emph{un} ADC tome cada muestra
en una fase distinta de períodos sucesivos. Como $F_s/f$ casi nunca es entero,
eso pasa solo; el algoritmo es:

\begin{enumerate}\itemsep1pt
  \item Capturar 2048 muestras a $\approx 2$\,MS/s, es decir, muchos períodos seguidos.
  \item Buscar todos los cruces ascendentes por el nivel medio, con histéresis,
        e interpolar entre muestras su posición fraccionaria $c_i$.
  \item Ajustar por mínimos cuadrados $c_i = t_0 + T\,i$. Con cientos de cruces,
        $T$ queda con un error menor a 1\,ppm. Primero se ajustan los cruces
        cercanos y la ventana se amplía de a poco, para numerarlos bien aunque
        se pierda alguno.
  \item Calcular la fase de cada muestra, $\varphi_k = \operatorname{frac}\big((k-t_0)/T\big)$,
        y acumularla en el casillero $\lfloor 256\,\varphi_k \rfloor$. Se guarda
        también dónde cayó dentro del casillero.
  \item Promediar cada casillero y llevar el valor a su \emph{centro} con la
        pendiente local, $y_c = \bar y - y'\,\bar\delta$. Sin esta corrección el
        error llega a 17\,LSB en los flancos; con ella baja a 0,4\,LSB.
  \item Interpolar los casilleros vacíos y enviar dos períodos (512 puntos) con
        $\Delta t_{eq} = T/(256\,F_s)$.
\end{enumerate}

\noindent\textbf{El micro se ``desafina'' a propósito.} Si $F_s/f$ queda cerca
de un cociente de enteros chicos, las muestras caen siempre en las mismas fases.
Por ejemplo, 500\,kHz a 2\,MS/s exactos da 4 muestras por período: solo 4 fases
distintas. Antes de reconstruir, el firmware prueba ARR\,=\,124\ldots131 (de 2,00
a 1,89\,MS/s), calcula cuántos de los 256 casilleros quedarían cubiertos y, si
otra $F_s$ cubre mejor, recaptura con esa. Es la versión en el tiempo del
desfase entre ADC.

\begin{figure}[h]
\centering
\begin{tikzpicture}
\begin{axis}[width=0.92\linewidth, height=6.3cm, xlabel={tiempo [ns]},
             ylabel={tensión [V]}, grid=major,
             legend style={font=\small, at={(0.5,1.03)}, anchor=south},
             legend columns=2, tick label style={font=\small},
             /pgf/number format/use comma, /pgf/number format/1000 sep={},
             ymin=0, ymax=3.3, xmin=0, xmax=2740]
  \addplot[only marks, mark=*, mark size=1.6pt, red] table[x=t_ns, y=v] {datos_ets_rt.dat};
  \addlegendentry{muestras reales (2,72 por período)}
  \addplot[thick, azul] table[x=t_ns, y=v] {datos_ets.dat};
  \addlegendentry{reconstrucción ETS (256 por período)}
\end{axis}
\end{tikzpicture}
\caption{Seno de 730\,kHz capturado a 1,98\,MS/s (simulación del ADC con ruido
de 1\,LSB y el código real de \texttt{osc\_core.c}). Los puntos rojos son las
muestras de los dos primeros períodos: a simple vista no se reconoce la forma.
La curva azul ordena por fase las 2048 muestras de la captura: equivale a
muestrear cada 5,35\,ns ($\approx 187$\,MS/s).}
\end{figure}

\begin{center}\small
\begin{tabular}{@{}rcccccc@{}}
\toprule
$f$ señal & Muestras/período & ARR elegido & Cobertura & Error de $f$ & Error de forma & Alineación \\
\midrule
50\,kHz  & 39,68 & 125 & 256/256 & 0,6\,ppm & 0,38\,LSB & 0,00\,ns \\
123,5\,kHz & 16,07 & 125 & 256/256 & 0,3\,ppm & 0,42\,LSB & 0,00\,ns \\
333\,kHz & 5,91 & 126 & 256/256 & 0,2\,ppm & 0,37\,LSB & $-0{,}23$\,ns \\
500\,kHz & 3,94 & 126 & 256/256 & 0,7\,ppm & 0,36\,LSB & $+0{,}23$\,ns \\
730\,kHz & 2,72 & 125 & 256/256 & 0,2\,ppm & 0,44\,LSB & 0,00\,ns \\
850\,kHz & 2,33 & 125 & 256/256 & 0,5\,ppm & 0,40\,LSB & $-1{,}24$\,ns \\
\bottomrule
\end{tabular}
\end{center}
Resultados de \texttt{make test} (seno de 2,4\,V\,pp con ruido de 1\,LSB). El
error de forma es RMS en LSB de 12 bits. El límite práctico es $f < F_s/2{,}2
\approx 900$\,kHz, porque hacen falta al menos dos muestras por período para
detectar los cruces, y la señal debe repetirse igual durante la captura. Para
eventos únicos se usa el modo en tiempo real. Por encima de unos pocos MHz
también limita el ancho de banda analógico de la entrada del ADC.

\subsection{Controles}

Mantener apretado el encoder FORMA cambia entre la página DDS y la del
osciloscopio, que arranca a adquirir solo. El generador sigue funcionando
mientras se mide: se puede ver la salida del DDS conectando PA4 a PA1.

\begin{center}
\begin{tabular}{@{}lll@{}}
\toprule
\textbf{Encoder} & \textbf{Acción} & \textbf{Página osciloscopio} \\
\midrule
FRECUENCIA & girar          & frecuencia de muestreo (horario = más rápido) \\
           & pulsar         & modo AUTO $\rightarrow$ NORMAL $\rightarrow$ ETS \\
           & mantener       & RUN / STOP \\
FORMA      & girar          & nivel de trigger (pasos de 33\,mV) \\
           & pulsar         & flanco ascendente / descendente \\
           & mantener       & volver a la página DDS \\
\bottomrule
\end{tabular}
\end{center}

\begin{figure}[h!]
\centering
\begin{tikzpicture}
  \node[inicio, minimum width=26mm] (lib) at (0,0) {LIBRE};
  \node[inicio, minimum width=26mm, fill=amarillo] (cap) at (8,0) {CAPTURANDO};
  \node[decision] (hay) at (8,-2.8) {¿hay trama para enviar?};
  \node[inicio, minimum width=26mm, fill=rojo] (env) at (8,-5.6) {ENVIANDO};
  \draw[flecha] (lib) -- node[etq, above, align=center]{RUN y UART libre:\\\texttt{osc\_hw\_start()}} (cap);
  \draw[flecha] (cap) -- node[etq, right, align=left]{DMA terminó: normalizar a 16 bits\\y procesar (tiempo real o ETS)} (hay);
  \draw[flecha] (hay) -- node[etq, right]{sí: encabezado + DMA de la UART} (env);
  \draw[flecha] (hay.west) -- node[etq, above, align=center]{no: NORMAL sin trigger,\\o ETS cambió de $F_s$} (0.5,-2.8) -- ([xshift=5mm]lib.south);
  \draw[flecha] (env.west) -- node[etq, below]{la UART terminó} (-0.5,-5.6) -- ([xshift=-5mm]lib.south);
\end{tikzpicture}
\caption{Máquina de estados de \texttt{osc\_task()}, en el lazo principal. Si
cambia la configuración, se aborta la captura y se vuelve a LIBRE.}
\end{figure}

\begin{figure}[p]
\centering
\begin{tikzpicture}[node distance=6mm and 10mm]
  \node[inicio] (s) {Captura terminada (modo ETS)};
  \node[ancho, below=of s] (a) {Mínimo, máximo, nivel medio e histéresis};
  \node[ancho, below=of a] (b) {Cruces ascendentes con histéresis:\\posición $c_i$ interpolada entre muestras};
  \node[ancho, below=of b] (c) {Regresión $c_i = t_0 + T\,i$ con ventana creciente\\$\Rightarrow$ período $T$ (en muestras) y fase inicial $t_0$};
  \node[decision, below=of c, text width=34mm] (d1) {¿periódica y al menos 4 períodos?};
  \node[procesoc, right=of d1] (rt) {mostrar en tiempo real};
  \node[ancho, below=10mm of d1] (e) {Para ARR = 124 \ldots 131: casilleros\\cubiertos por las fases $\operatorname{frac}(k/T')$};
  \node[decision, below=of e, text width=34mm] (d2) {¿otra $F_s$ cubre mejor?};
  \node[procesoc, right=of d2] (re) {cambiar ARR\\y recapturar};
  \node[ancho, below=10mm of d2] (f) {Para cada muestra $k$:\\$\varphi = \operatorname{frac}((k - t_0)/T)$, \; $b = \lfloor 256\,\varphi \rfloor$\\acumular valor y desvío dentro del casillero};
  \node[ancho, below=of f] (g) {Promediar y corregir al centro: $y_c = \bar y - y'\,\bar\delta$};
  \node[ancho, below=of g] (h) {Interpolar casilleros vacíos};
  \node[io, below=of h] (i) {Trama: 2 períodos $\times$ 256 puntos\\$\Delta t = T/(256\,F_s)$, $f = F_s/T$};
  \draw[flecha] (s) -- (a);
  \draw[flecha] (a) -- (b);
  \draw[flecha] (b) -- (c);
  \draw[flecha] (c) -- (d1);
  \draw[flecha] (d1) -- node[etq, above]{no} (rt);
  \draw[flecha] (d1) -- node[etq, right]{sí} (e);
  \draw[flecha] (e) -- (d2);
  \draw[flecha] (d2) -- node[etq, above]{sí} (re);
  \draw[flecha] (d2) -- node[etq, right]{no} (f);
  \draw[flecha] (f) -- (g);
  \draw[flecha] (g) -- (h);
  \draw[flecha] (h) -- (i);
\end{tikzpicture}
\caption{Procesamiento en tiempo equivalente (\texttt{process\_ets()} y \texttt{osc\_core.c}).}
\end{figure}

\subsection{Acondicionamiento de la entrada}

La entrada va directo al ADC: el rango es de 0 a 3,3\,V
($V_{REF+} = V_{DDA}$ en el LQFP48). Recomendaciones:
\begin{itemize}\itemsep0pt
  \item \textbf{Protección:} resistencia serie de 1\,k$\Omega$ y diodos Schottky
        a 3,3\,V y a GND (por ejemplo, BAT54S).
  \item \textbf{Buffer:} a 2\,MS/s el ADC muestrea durante solo 80\,ns
        (2,5 ciclos), así que la fuente tiene que ser de baja impedancia. Un
        operacional \emph{rail-to-rail} rápido como seguidor, por ejemplo un
        OPA350 o similar, con un capacitor de 1--10\,nF a la salida para el
        pico de carga del ADC. En las bases lentas el firmware alarga solo el
        tiempo de muestreo.
  \item \textbf{Señales bipolares:} divisor resistivo más un \emph{offset} de
        1,65\,V con el mismo operacional.
\end{itemize}

% =========================================================================
\clearpage
\section{Conexionado (LQFP48)}

\begin{center}
\begin{tabular}{@{}llll@{}}
\toprule
\textbf{Pin} & \textbf{Función} & \textbf{AF} & \textbf{Conexión} \\
\midrule
PA4  & DAC1\_OUT1 & analóg. & Salida CH1 $\rightarrow$ filtro $\rightarrow$ BNC \\
PA5  & DAC1\_OUT2 & analóg. & Salida CH2 $\rightarrow$ filtro $\rightarrow$ BNC \\
PA8 / PA9 & TIM1\_CH1 / CH2 & AF1 & Encoder FRECUENCIA (A / B) \\
PA10 & GPIO entrada, pull-up & -- & Pulsador FRECUENCIA (cambia el paso) \\
PA6 / PA7 & TIM3\_CH1 / CH2 & AF2 & Encoder FORMA (A / B) \\
PB0  & GPIO entrada, pull-up & -- & Pulsador FORMA (modo de CH2) \\
PA1  & ADC1\_INP1 & analóg. & \textbf{Entrada del osciloscopio} (0 a 3,3\,V) \\
PA2  & USART2\_TX & AF7 & Estado + tramas del osciloscopio, 2\,Mbaud 8N1 \\
PA13 / PA14 & SWDIO / SWCLK & -- & ST-LINK \\
\midrule
PB6 / PB7 & I2C1 SCL / SDA & AF4 & \emph{Reservado:} pantalla I$^2$C (SSD1306) \\
PB13/15/12 & SPI2 SCK/MOSI/CS & AF5 & \emph{Reservado:} pantalla SPI (ST7735, ILI9341) \\
PB14 / PB10 & GPIO & -- & \emph{Reservado:} D/C y RESET de la pantalla \\
\bottomrule
\end{tabular}
\end{center}

\noindent\textbf{Recomendaciones de hardware.} Encoders con el común a GND y una
red RC de 10\,k$\Omega$/10\,nF en A y B. En el LQFP48, $V_{REF+}$ está unido a
$V_{DDA}$: alimentá $V_{DDA}$ con un filtro LC aparte para bajar el ruido del DAC.
Con el buffer interno activado la salida no llega a los rieles (unos 0,2\,V de
cada lado). Para amplitud y offset ajustables agregá una etapa con un
amplificador operacional después del filtro reconstructor.

\noindent\textbf{Pantalla futura.} Hay tres funciones \texttt{weak} (vacías):
\texttt{display\_init()}, \texttt{display\_update(estado)} y
\texttt{display\_scope(trama, puntos)}. Basta con
crear \texttt{display.c} con esas funciones. La primera recibe la página activa,
la frecuencia pedida y la real, el paso (para dibujar el cursor), la forma de
onda, el modo de CH2 y la configuración del osciloscopio. La segunda recibe
cada trama lista para dibujar: puntos de 16 bits, $\Delta t$ entre puntos y
mediciones. Ni el motor DDS ni la adquisición se tocan.

\end{document}
